% Lecture notes: Week 11, Lecture 2 -- Strain Transformations
% To hide this lecture from the website, add a line containing only:  % draft
% To set the date shown on the website, add a line like:  % posted: 2026-09-02
\documentclass[11pt]{article}
\usepackage{mech2030notes}
\begin{document}

% ##################################################################
%  WEEK 11, LECTURE 2
% ##################################################################
\lecture{11}{2}{Strain Transformations}

\noindent
Just like we have 2D \emph{stress} states which combine shear and normal stresses, we can
also have 2D \emph{strain} states which combine shear and normal strains.

\begin{center}
\begin{minipage}[t]{0.45\textwidth}
\centering
\textbf{2D stress state} (plane stress)\\[0.6em]
\begin{tikzpicture}[scale=1.1]
  \draw[thick, fill=gray!12] (-0.7,-0.7) rectangle (0.7,0.7);
  \draw[axis] (0,0) -- (0.45,0) node[right] {\scriptsize$x$};
  \draw[axis] (0,0) -- (0,0.45) node[left] {\scriptsize$y$};
  \draw[force] (0.75,0) -- (1.4,0) node[right] {$\sigma_x$};
  \draw[force] (-0.75,0) -- (-1.4,0) node[left] {$\sigma_x$};
  \draw[force] (0,0.75) -- (0,1.4) node[above] {$\sigma_y$};
  \draw[force] (0,-0.75) -- (0,-1.4) node[below] {$\sigma_y$};
  \draw[force] (0.9,-0.45) -- (0.9,0.45);
  \draw[force] (-0.9,0.45) -- (-0.9,-0.45);
  \draw[force] (-0.45,0.9) -- (0.45,0.9) node[right] {$\tau_{xy}$};
  \draw[force] (0.45,-0.9) -- (-0.45,-0.9) node[left] {$\tau_{xy}$};
\end{tikzpicture}
\end{minipage}%
\hfill
\begin{minipage}[t]{0.45\textwidth}
\centering
\textbf{2D strain state} (plane strain)\\[0.6em]
\begin{tikzpicture}[scale=1.1]
  \draw[thick, fill=gray!12] (-0.7,-0.7) rectangle (0.7,0.7);
  \draw[axis] (0,0) -- (0.45,0) node[right] {\scriptsize$x$};
  \draw[axis] (0,0) -- (0,0.45) node[left] {\scriptsize$y$};
  \draw[force] (0.75,0) -- (1.4,0) node[right] {$\epsilon_x$};
  \draw[force] (-0.75,0) -- (-1.4,0) node[left] {$\epsilon_x$};
  \draw[force] (0,0.75) -- (0,1.4) node[above] {$\epsilon_y$};
  \draw[force] (0,-0.75) -- (0,-1.4) node[below] {$\epsilon_y$};
  \draw[force] (0.9,-0.45) -- (0.9,0.45);
  \draw[force] (-0.9,0.45) -- (-0.9,-0.45);
  \draw[force] (-0.45,0.9) -- (0.45,0.9) node[right] {$\gamma_{xy}$};
  \draw[force] (0.45,-0.9) -- (-0.45,-0.9) node[left] {$\gamma_{xy}$};
\end{tikzpicture}
\end{minipage}
\end{center}

% ==================================================================
\section*{Strain Transformation Equations}
% ==================================================================

\noindent
The 2D strain transformation relations are \emph{almost} the same as the 2D stress
transformation relations (Week~11, Lecture~1), with \textbf{one important difference}.
The difference is just due to the particular definition of the shear strain we use.

\begin{center}
\renewcommand{\arraystretch}{1.6}
\textbf{Substitute:}\qquad
\begin{tabular}{c @{\quad$\Longleftrightarrow$\quad} c}
  $\sigma_{x'}$ & $\epsilon_{x'}$ \\
  $\sigma_{y'}$ & $\epsilon_{y'}$ \\
  \multicolumn{2}{c}{\fcolorbox{red!70!black}{red!8}{%
    $\tau_{x'y'} \quad\Longleftrightarrow\quad \dfrac{\gamma_{x'y'}}{2}$}}
\end{tabular}
\end{center}

\noindent
\begin{minipage}[t]{0.48\textwidth}
\centering
\textbf{General 2D stress transformation equations:}
\begin{align*}
  \sigma_{x'} &= \frac{\sigma_x + \sigma_y}{2} + \frac{\sigma_x - \sigma_y}{2}\cos 2\theta
                 + \tau_{xy}\sin 2\theta \\
  \sigma_{y'} &= \frac{\sigma_x + \sigma_y}{2} - \frac{\sigma_x - \sigma_y}{2}\cos 2\theta
                 - \tau_{xy}\sin 2\theta \\
  \tau_{x'y'} &= -\frac{\sigma_x - \sigma_y}{2}\sin 2\theta + \tau_{xy}\cos 2\theta
\end{align*}
\end{minipage}%
\hfill
\begin{minipage}[t]{0.48\textwidth}
\centering
\textbf{General 2D strain transformation equations:}\\[0.4em]
\fbox{\fbox{\begin{minipage}{0.93\textwidth}
\vspace{-0.6em}
\begin{align*}
  \epsilon_{x'} &= \frac{\epsilon_x + \epsilon_y}{2} + \frac{\epsilon_x - \epsilon_y}{2}\cos 2\theta
                   + \frac{\gamma_{xy}}{2}\sin 2\theta \\
  \epsilon_{y'} &= \frac{\epsilon_x + \epsilon_y}{2} - \frac{\epsilon_x - \epsilon_y}{2}\cos 2\theta
                   - \frac{\gamma_{xy}}{2}\sin 2\theta \\
  \frac{\gamma_{x'y'}}{2} &= -\frac{\epsilon_x - \epsilon_y}{2}\sin 2\theta
                             + \frac{\gamma_{xy}}{2}\cos 2\theta
\end{align*}
\end{minipage}}}
\end{minipage}

% ==================================================================
\section*{Mohr's Circle for Strains}
% ==================================================================

\noindent
This means \emph{you can also draw Mohr's circle for strains} instead of stresses, using the
same substitutions. The same substitutions also apply to the expressions for the principal
stresses/strains and the maximum shear:

\noindent
\begin{minipage}[t]{0.48\textwidth}
\centering
\textbf{Stresses:}
\begin{align*}
  \sigma_{\text{avg}} &= \frac{\sigma_x + \sigma_y}{2} \\
  R = \tau_{\max} &= \sqrt{\left(\frac{\sigma_x - \sigma_y}{2}\right)^2 + \tau_{xy}^2} \\
  \sigma_1 &= \sigma_{\text{avg}} + R \\
  \sigma_2 &= \sigma_{\text{avg}} - R
\end{align*}
\end{minipage}%
\hfill
\begin{minipage}[t]{0.48\textwidth}
\centering
\textbf{Strains:}\\[0.4em]
\fbox{\fbox{\begin{minipage}{0.93\textwidth}
\vspace{-0.6em}
\begin{align*}
  \epsilon_{\text{avg}} &= \frac{\epsilon_x + \epsilon_y}{2} \\
  R = \frac{\gamma_{\max}}{2} &= \sqrt{\left(\frac{\epsilon_x - \epsilon_y}{2}\right)^2
                                  + \left(\frac{\gamma_{xy}}{2}\right)^2} \\
  \epsilon_1 &= \epsilon_{\text{avg}} + R \\
  \epsilon_2 &= \epsilon_{\text{avg}} - R
\end{align*}
\end{minipage}}}
\end{minipage}

% ==================================================================
\section*{Strain Gauge Rosettes}
% ==================================================================

\noindent
\begin{minipage}[c]{0.55\textwidth}
The 2D strain transformations have an important practical use in measuring shear strains.
\begin{itemize}
  \item Strain gauges can measure \emph{normal} strains directly, but \emph{not} shear
        strains.
  \item Shear strains are instead computed using a \emph{strain gauge rosette}: several
        gauges bonded at a point, each oriented in a different direction.
\end{itemize}
For the 3-element rosette shown, clearly we have
\[
  \epsilon_a = \epsilon_x, \qquad \epsilon_b = \epsilon_y.
\]
\begin{itemize}[label=$\rightarrow$]
  \item How do we calculate $\gamma_{xy}$?
\end{itemize}
\end{minipage}%
\hfill
\begin{minipage}[c]{0.42\textwidth}
\centering
\begin{tikzpicture}
  \draw[axis] (0,0) -- (3.0,0) node[right] {$x$};
  \draw[axis] (0,0) -- (0,3.0) node[above] {$y$};
  \fill (0,0) circle (0.06);
  \draw[thin, dashed] (0,0) -- (45:3.0);
  % gauge 1 along x
  \draw[thick, fill=white] (1.0,-0.09) rectangle (1.9,0.09);
  \node[below] at (1.45,-0.1) {$\epsilon_a$};
  % gauge 2 along y
  \draw[thick, fill=white] (-0.09,1.0) rectangle (0.09,1.9);
  \node[left] at (-0.1,1.45) {$\epsilon_b$};
  % gauge 3 at 45 deg
  \begin{scope}[rotate=45]
    \draw[thick, fill=white] (1.1,-0.09) rectangle (2.0,0.09);
  \end{scope}
  \node[above left] at (45:1.6) {$\epsilon_c$};
  \draw[thinarrow] (2.45,0) arc (0:45:2.45);
  \node at (22:2.75) {$\ang{45}$};
\end{tikzpicture}
\end{minipage}

% ==================================================================
\section*{Example}
% ==================================================================

\noindent
\textbf{Measurement:} The rosette above measures
\[
  \left\{
  \begin{aligned}
    \epsilon_a &= 100\,\mu\epsilon \quad (= 100 \times 10^{-6}) \\
    \epsilon_b &= 900\,\mu\epsilon \\
    \epsilon_c &= 200\,\mu\epsilon
  \end{aligned}
  \right.
\]
\textbf{Q:} What is $\gamma_{xy}$?

\bigskip
\noindent
$\rightarrow$ Clearly,
\[
  \left\{
  \begin{aligned}
    \epsilon_x &= \epsilon_a = 100\,\mu\epsilon \\
    \epsilon_y &= \epsilon_b = 900\,\mu\epsilon
  \end{aligned}
  \right.
\]

\noindent
$\rightarrow$ Then, $\epsilon_c$ is $\epsilon_{x'}$ at $\theta = \ang{45}$:
\[
  \epsilon_c = \epsilon_{x'}(\theta = \ang{45})
  = \frac{\epsilon_a + \epsilon_b}{2}
    + \frac{\epsilon_a - \epsilon_b}{2}\underbrace{\cos(\ang{90})}_{=\,0}
    + \frac{\gamma_{xy}}{2}\underbrace{\sin(\ang{90})}_{=\,1}
\]
\[
  \hookrightarrow\quad
  \epsilon_c = \frac{\epsilon_a + \epsilon_b}{2} + \frac{\gamma_{xy}}{2}
\]
Solving for $\gamma_{xy}$:
\begin{align*}
  \gamma_{xy} &= 2\left(\epsilon_c - \frac{\epsilon_a + \epsilon_b}{2}\right) \\
  \gamma_{xy} &= 2\left(200 - \frac{100 + 900}{2}\right)\mu\epsilon
\end{align*}
\[
  \fbox{\fbox{$\gamma_{xy} = -600\,\mu\epsilon$}}
\]

\end{document}
